\section{Limitations and threats to validity}
\label{sec:limits}

\paragraph{The ground truth covers one effect on one data set}
C7 measures the response of a store's category units to a sustained change in
its regular prices, on Dominick's. It says nothing directly about promotions,
displays or temporary cuts, or about M5 and Frat, where no such variation is used. A model can be right about regular prices and wrong about
promotions; the known-answer study shows how far the other checks go in
covering the rest.

\paragraph{The deviations are inferred}
The public files do not mark the experimental stores or weeks. Some deviations
may be zone moves made for other reasons, and a store moved for a reason that
also moves its sales would bias $\tau_e$. The event study and the placebo
address this for trends that start before the price moves; they cannot rule out
a shock that starts in the same week.

\paragraph{Few control stores}
Zones have few stores, so each scored episode has \zControlsLo{} to
\zControlsHi{} controls. Intervals clustered
by store with so few clusters are wide and only approximately calibrated; the
pooled slope across episodes carries the comparison.

\paragraph{The baseline is a choice}
Every lift is measured against nothing planned, with the regular price taken
as the highest of the last 12 periods. A different rule gives different
baselines and lifts. The rule is the models' own and is applied the same way to
every model and to the truth in the known-answer study.

\paragraph{The versions are not a designed sample}
The audited versions are what the companion paper saved. They share data and
differ in many settings at once, so the correlation across versions is
descriptive; it is computed within data set and tested by permutation within
data set, but versions are not independent draws, and the windows of one
version differ in difficulty, so the result depends on whether version-windows
or versions are the unit. Both are reported.

\paragraph{The audit's definition of nothing planned}
The audit built ``nothing planned'' from each run's configuration. Using the
regular price the model stores did not remove the C1 violations in the
fourteen rerun runs, and configurations that read a recorded regular price or
contain new items were not tested further. We did not find the cause, so C1
failures on real data may come from the audit's definition as much as from the
models.

\paragraph{Refitted versions}
The Dominick's final world models are refits with unchanged settings, not the
runs whose scores the companion paper reports; seed-to-seed differences between
the two are expected and are measured by C3.

\paragraph{The generated world is simpler than retail}
The known-answer data are multiplicative in the plan, with one shock process
and no cannibalisation between items. They test whether a check can catch a
failure, not how often the failure occurs in practice.
